% Einheit 4 – Verdopplungs- und Halbwertszeit
\setcounter{aufgabe}{35}
\hypertarget{e4}{}\einheitenkopf*{Einheit 4 von 5 · Verdopplungs- und Halbwertszeit}
\zweigzeile{Hier lernst du, die Zeit zu bestimmen, nach der sich ein Wert verdoppelt oder halbiert hat \,·\, neu in diesem Jahr \,·\, P10 \,·\, baut auf: Einheit 3, Wertetabelle und Graph lesen, Potenz mit dem Taschenrechner}

\begin{aufgabe}{Ich erkenne, ob in einer Frage die Zeit oder der Bestand verdoppelt oder halbiert wird. Kreuze an. Du musst nichts ablesen.}
\begin{teile}
\teil Nach wie vielen Jahren hat sich das Guthaben verdoppelt? \hfill \kreuz{Zeit}\kreuz{Bestand}
\teil Wie viel Geld ist auf dem Konto, wenn man doppelt so lange spart? \hfill \kreuz{Zeit}\kreuz{Bestand}
\teil Wann ist nur noch die Hälfte des Wirkstoffs im Blut? \hfill \kreuz{Zeit}\kreuz{Bestand}
\teil Wie viele Bakterien sind es nach der halben Zeit? \hfill \kreuz{Zeit}\kreuz{Bestand}
\end{teile}
\end{aufgabe}

\begin{aufgabe}{Ich kann die Verdopplungszeit und die Halbwertszeit bestimmen. Bilde den Zielwert und kreuze an, welche Zeit gesucht ist.}
\beispiel{\text{Anfangswert } 350\,\text{g, der Bestand nimmt ab} &\;\rightarrow\; \text{Zielwert } 175\,\text{g, gesucht ist die Halbwertszeit}}
\begin{teile}
\teil Anfangswert $90$, der Bestand nimmt zu. \hfill Zielwert: \leerfeld \kreuz{Verdopplungszeit}\kreuz{Halbwertszeit}
\teil Anfangswert $1\,400$, der Bestand nimmt ab. \hfill Zielwert: \leerfeld \kreuz{Verdopplungszeit}\kreuz{Halbwertszeit}
\teil Anfangswert $2{,}5\,\text{mg}$, der Bestand nimmt ab. \hfill Zielwert: \leerfeld \kreuz{Verdopplungszeit}\kreuz{Halbwertszeit}
\teil Anfangswert $725$, der Bestand nimmt zu. \hfill Zielwert: \leerfeld \kreuz{Verdopplungszeit}\kreuz{Halbwertszeit}
\end{teile}
Lies die gesuchte Zeit aus der Tabelle ab.
\begin{teile}
\teil Ein Stoff zerfällt. \hfill Halbwertszeit: \leerfeld[Tage]
\end{teile}
\sachtabelle{lccccc}{Zeit in Tagen & 0 & 4 & 8 & 12 & 16}{Menge in g & 1\,600 & 1\,131 & 800 & 566 & 400}
\begin{teile}
\teil Ein Bestand wächst. Der doppelte Wert steht nicht genau in der Tabelle. Gib die Zeit an, die am nächsten liegt. \hfill Verdopplungszeit: etwa \leerfeld[Jahre]
\end{teile}
\sachtabelle{lcccccccc}{Jahr & 0 & 1 & 2 & 3 & 4 & 5 & 6 & 7}{Bestand & 300 & 339 & 383 & 433 & 489 & 553 & 625 & 706}
\end{aufgabe}

\begin{aufgabe}{Ich kann die Verdopplungszeit und die Halbwertszeit bestimmen – weiter. Ein Wirkstoff wird im Körper abgebaut. Der Graph zeigt die Menge in mg, $t$ ist die Zeit in Stunden.}

\begin{ksys}[xmin=0,xmax=10,ymin=0,ymax=360,ystep=40,ablesen,xlabel=$t$ in h,ylabel=mg]
\funktion[7]{320*0.5^(\x/3)}{}
\end{ksys}

\begin{teile}
\teil Lies die Halbwertszeit am Graphen ab. \quad \leerfeld[h]
\teil Beschreibe in zwei Sätzen, wie du vorgegangen bist.
\end{teile}
\end{aufgabe}

\begin{aufgabe}{Ich kann die Verdopplungszeit und die Halbwertszeit bestimmen – weiter. Rechne mit dem Wachstumsfaktor.}
\begin{teile}
\teil Ein Bestand von $40$ Tieren wächst jedes Jahr um $15\,\%$. Rechne Schritt für Schritt, bis sich der Bestand verdoppelt hat. Nach wie vielen Jahren ist das erstmals der Fall? \quad \leerfeld[Jahre]
\teil Zeige: Auch ein Bestand von $300$ Tieren braucht bei $15\,\%$ Wachstum genauso viele Jahre, bis er sich verdoppelt hat.
\teil Mit dem Taschenrechner geht es genau: Die Verdopplungszeit bei $15\,\%$ ist $t = \log(2) : \log(1{,}15)$; $\log$ ist eine Taste deines Taschenrechners. Berechne $t$ und runde auf eine Nachkommastelle. \quad \leerfeld[Jahre]
\teil Der Graph zeigt die Zahl der Elektroautos in einer Stadt (in Tausend), $t$ in Jahren seit Beginn der Zählung. Bestimme mithilfe des Graphen, nach wie vielen Jahren sich die Zahl verdoppelt hat. Beschreibe dein Vorgehen. (P10 2020)
\end{teile}

\begin{ksys}[xmin=0,xmax=12,ymin=0,ymax=4.8,ystep=0.4,ablesen,xlabel=$t$ in Jahren,ylabel=Tausend]
\funktion[10]{1.6*1.09^\x}{}
\end{ksys}

\end{aufgabe}

\begin{aufgabe}{Ich finde den Fehler beim Bestimmen der Halbwertszeit. Ein Stoff zerfällt:}
\sachtabelle{lccccc}{Zeit in Stunden & 0 & 2 & 4 & 6 & 8}{Menge in mg & 500 & 420 & 354 & 297 & 250}
Emil schreibt: „Die Tabelle geht bis 8 Stunden. Die Hälfte davon sind 4 Stunden, also ist die Halbwertszeit 4 Stunden.“
\begin{teile}
\teil Finde den Fehler, benenne ihn und gib die richtige Halbwertszeit an.
\teil Bestimme selbst die Halbwertszeit. \hfill Halbwertszeit: \leerfeld[Jahre]
\end{teile}
\sachtabelle{lccccc}{Zeit in Jahren & 0 & 9 & 18 & 27 & 36}{Menge in g & 2\,400 & 1\,697 & 1\,200 & 849 & 600}
\end{aufgabe}

\begin{aufgabe}{Ich kann begründen, warum die Verdopplungszeit immer gleich lang ist. Entscheide ohne Rechnung und begründe.}
\begin{teile}
\teil Ein Bestand verdoppelt sich alle $7$ Jahre. Wie viel Mal so groß wie am Anfang ist er nach $21$ Jahren?
\teil Warum genügt es, den doppelten Wert einmal zu suchen? Warum dauert die nächste Verdopplung genauso lang?
\teil Ein Bestand wächst jedes Jahr um denselben Betrag. Hat er auch eine feste Verdopplungszeit?
\end{teile}
\end{aufgabe}

\begin{aufgabe}{Ich kann mit der Halbwertszeit eine Entscheidung treffen. Ein Pflanzenschutzmittel zerfällt im Boden mit einer Halbwertszeit von $12$ Tagen. Auf ein Feld werden $8\,\text{kg}$ ausgebracht. Gemüse darf erst gesät werden, wenn weniger als $1\,\text{kg}$ im Boden ist.}
\begin{teile}
\teil Wie viel Mittel ist nach $12$, nach $24$ und nach $36$ Tagen noch im Boden?
\teil Reicht es, vier Wochen ($28$ Tage) zu warten? Begründe.
\teil Ab wann darf gesät werden?
\end{teile}
\end{aufgabe}
